% =====================================================================
% Springer Nature sn-jnl submission version
% Paper: Relational Time-of-Arrival from an Event Current and its Continuum Clock Limit
% Author: Omar Iskandarani
% Notes:
% - Single .tex file (no \input)
% =====================================================================

% \documentclass[referee,pdflatex,sn-mathphys-num]{sn-jnl} % optional for review
\documentclass[pdflatex,sn-mathphys-num]{sn-jnl}

% =========================
% Core packages (lean)
% =========================
\usepackage[T1]{fontenc}
\usepackage[utf8]{inputenc}
\usepackage{lmodern}
\usepackage{microtype}
\usepackage{graphicx}
\usepackage{booktabs}
\usepackage{amsmath,amssymb,amsfonts}
\usepackage{amsthm}
\usepackage{mathtools}
\usepackage{bm}
\usepackage{bbm}

% =========================
% Theorem environments
% =========================
\theoremstyle{thmstylethree}
\newtheorem{definition}{Definition}[section]
\newtheorem{lemma}{Lemma}[section]

\theoremstyle{thmstyleone}
\newtheorem{proposition}{Proposition}[section]

\theoremstyle{thmstyletwo}
\newtheorem{remark}{Remark}[section]

\numberwithin{equation}{section}

% =========================
% Macros
% =========================
\newcommand{\dd}{\mathrm{d}}
\newcommand{\Sig}{\Sigma}
\newcommand{\Jev}{j^\mu_{\mathrm{ev}}}
\newcommand{\Nev}{N_{\mathrm{ev}}}
\newcommand{\F}{\mathcal{F}}
\newcommand{\cO}{\mathcal{O}}
\newcommand{\eps}{\varepsilon}

% QFT macros
\newcommand{\W}{W}
\newcommand{\Gret}{G_{\mathrm{ret}}}
\newcommand{\Gplus}{G^{+}}
\newcommand{\Kzero}{K_{0}}

% =========================
% Metadata
% =========================
\newcommand{\papertitle}{Relational Time-of-Arrival as a Covariant Field Observable: Integrable Event Currents, Clock Holonomy, and a Continuum Clock Limit}
\newcommand{\papershorttitle}{Relational Time-of-Arrival}
\newcommand{\paperemail}{info@omariskandarani.com}
\newcommand{\paperlocation}{Groningen}
\newcommand{\paperkeywords}{relational time, time of arrival, quantum time observables, covariant observables, event currents, clock holonomy, quantum field theory}
\newcommand{\paperdoi}{10.5281/zenodo.20486634}
\newcommand{\paperorcid}{0009-0006-1686-3961}

\begin{document}
    
    \title[\papershorttitle]{\papertitle}
    
    \author*[1]{\fnm{Omar} \sur{Iskandarani}}\email{\paperemail}
    
    \affil*[1]{\orgname{Independent Researcher},
        \orgaddress{\city{\paperlocation}, \country{The Netherlands}}
    }
    
    \abstract{%
        We formulate quantum time-of-arrival (TOA) as a covariant relational field observable rather than as a self-adjoint time operator. Detector crossings are described by a conserved matter current $J^\mu$ through a world-tube $\Sigma$, while clock readings are generated by an effective event-count current $j^\mu_{\mathrm{ev}}$. The construction separates conservation from integrability: conservation stabilizes event counts, but a scalar continuum clock $T(x)$ exists only when the coarse-grained event-current one-form is locally integrable. The obstruction is the event-current vorticity
        $\Omega^{\mathrm{ev}}_{\mu\nu}=\partial_\mu(j^{\mathrm{ev}}_\nu)_\ell-\partial_\nu(j^{\mathrm{ev}}_\mu)_\ell$.
        Integrable currents yield a flux-conditioned TOA density; non-integrable currents yield path-conditioned TOA distributions governed by clock holonomy. In a $1+1$-dimensional Klein--Gordon example with a gapped quadratic clock EFT, Gaussian detector smearing gives the explicit regulated clock variance
        $\sigma_{\tau,\ell}^2=(4\pi)^{-1}K_0(\mu_\tau^2\ell^2)$.
        This provides a detector-level framework in which TOA statistics depend explicitly on flux, geometry, clock-current admissibility, holonomy, and clock-sector correlators.
    }
    
    \keywords{\paperkeywords}
    
    \maketitle

% ======================================================================
    \section*{Notation summary}
% ======================================================================
        
        \begin{center}
            \begin{tabular}{ll}
                $J^\mu$ & conserved matter current (detector flux) \\
                $j^\mu_{\mathrm{ev}}$ & conserved event-count current (clock sector) \\
                $N_{\mathrm{ev}}$ & discrete event-count clock \\
                $T(x)$ & continuum clock field (IR limit, when integrability holds) \\
                $\Sigma$ & detector world-tube \\
                $\mathcal{F}$ & detector flux functional $n_\mu J^\mu$ \\
                $\Omega^{\mathrm{ev}}_{\mu\nu}$ & event-current vorticity / clock-integrability obstruction \\
                $(\cdot)_\ell$ & coarse-graining over spacetime scale $\ell$
            \end{tabular}
        \end{center}

% ======================================================================
    \section{Introduction}
% ======================================================================
        
        Time is usually supplied to quantum theory as an external parameter. Operationally, however, time is read from physical degrees of freedom, and in relativistic settings those degrees of freedom are local. This mismatch is visible already in nonrelativistic quantum mechanics: Pauli-type objections constrain the existence of a self-adjoint time operator canonically conjugate to a semibounded Hamiltonian in generic settings \cite{Pauli1980,Muga2008}. In quantum field theory the situation is sharper, because localization, causality, and detector modelling enter from the start \cite{Muga2008,ColosiOeckl2025TimeOperatorTOA}.
        
        Many successful TOA constructions therefore adopt explicitly operational definitions: POVM-based prescriptions such as Kijowski's \cite{Kijowski1974} and later developments \cite{Giannitrapani1997,Muga2008}, as well as detector models and measurement couplings \cite{ColosiOeckl2025TimeOperatorTOA}. These approaches produce useful arrival-time statistics, but typically tie the observable to a chosen clock variable, a preferred slicing, or detector microphysics.
        
        In this paper we take a different starting point. We define TOA as a \emph{relational field observable}: detector crossings are binned by a \emph{physical clock} constructed from an event-count current. Concretely, we introduce (i) a conserved matter current $J^\mu$ that determines flux through a detector world-tube $\Sigma$, and (ii) an effective current $j^\mu_{\mathrm{ev}}$ whose integral counts localized detector-relevant events and thereby defines a discrete clock. The passage from this discrete clock to a continuum scalar clock is not automatic. Conservation of $j^\mu_{\mathrm{ev}}$ controls deformation-invariance of counts, while scalar-clock existence requires the additional local integrability condition $\partial_{[\mu}(j^{\mathrm{ev}}_{\nu]})_\ell=0$ after coarse-graining. If this condition fails, the framework does not break down; instead it predicts a path-dependent clock functional with a holonomy determined by the event-current vorticity.
        
        The payoff is threefold. First, the TOA observable is manifestly covariant and does not require postulating a self-adjoint time operator. Second, the construction cleanly separates event-count conservation from clock-field integrability, making the obstruction to continuum time explicit. Third, once an admissible clock sector is specified, its correlators enter arrival-time statistics in a parameter-controlled way.
        
        \paragraph{Relation to existing time-of-arrival formalisms.}
            The TOA literature offers several workable definitions---POVM prescriptions, explicit detector models, and path-integral or histories-based constructions. Our goal is not to compete with their operational success, but to isolate a covariant object that can be stated directly at the field level. The key move is to treat the clock as part of the physics: TOA is defined by conditioning flux through $\Sigma$ on a clock reading built from an event current. In that sense, detector dependence is not removed; it is made explicit and parameterized by the world-tube geometry, the flux current, the admissible event current, and the clock-sector correlators.
        
        \paragraph{Relation to Page--Wootters, partial observables, and relational-clock frameworks.}
            The present construction is complementary to Page--Wootters-type and relational-clock frameworks rather than a replacement for them \cite{PageWootters1983,Wootters1984}. In those approaches, time is typically recovered through conditional states or correlations between a clock subsystem and the rest of the system. The present work is also close in spirit to the partial/complete-observable programme in generally covariant systems, where physical predictions are formulated as correlations between dynamical quantities rather than as functions of an external background time \cite{Rovelli2002PartialObservables,Dittrich2007PartialComplete,Tambornino2012RelationalObservables,Marolf1995AlmostIdealClocks}. Here, by contrast, the emphasis is not on Hilbert-space factorization or canonical Dirac observables, but on a covariant field-level detector functional: the flux through a world-tube $\Sigma$ conditioned on a local clock readout generated by an effective event current. The novelty of the present formulation is therefore structural rather than metaphysical: it packages relational time-of-arrival directly as a covariant detector functional, separates conservation from integrability in the event-current clock, and identifies clock holonomy as the obstruction to a single-valued continuum time variable. In the non-integrable sector this obstruction is not merely diagnostic: it enters a path-conditioned TOA distribution, so that different admissible clock-path families can lead to experimentally distinguishable readout statistics. In this sense the construction is in dialogue with related perspectives on relational clocks, quantum reference frames, and operational notions of time \cite{GambiniPullin2021,HohnSmith2020,Giovannetti2015}, while keeping the main object of interest at the level of a covariant detector functional rather than a conditional state construction.
            
            Rather than revising measurement postulates or postulating a generalized time operator, we keep the standard field-theoretic setting and change only what is taken to be primitive: TOA is defined as a relational field-functional built from conserved currents and a physical clock sector. The absence of a universal time operator is then not a gap to be patched, but a cue that TOA should be formulated relationally through correlations between physical degrees of freedom.

        \paragraph{Scope and assumptions.}
            Throughout this work, the event current $j^\mu_{\mathrm{ev}}$ is treated as a \emph{physical but model-agnostic} structure. It may be fundamental (e.g.\ a primitive detection or interaction current) or emergent (e.g.\ a coarse-grained description of localized spacetime processes such as interaction vertices, defect crossings, or excitation worldlines). The approximate conservation $\nabla_\mu j^\mu_{\mathrm{ev}}\approx 0$ is expected to hold whenever the relevant events are robust under smooth deformations on the scales of interest, with violations suppressed by the ultraviolet resolution scale. The existence of a continuum scalar clock requires the independent integrability condition on the coarse-grained one-form $(j^{\mathrm{ev}}_\mu)_\ell$; without it one obtains a path-dependent clock rather than a global clock field.
            
            The coarse-graining map $(\cdot)_\ell$ is assumed to be a local, causal smoothing operation over spacetime regions of characteristic size $\ell$, preserving current conservation in the infrared. No specific kernel is fixed; admissible choices include covariant averaging, block-spin--type maps, or detector-induced temporal binning, provided they respect locality and current conservation at scales $\gg\ell$.
            Admissibility criteria for $j^\mu_{\mathrm{ev}}$ in a given experimental context are discussed further in Section~\ref{sec:discussion}.

% ======================================================================
    \section{Detector World-Tubes and Matter Flux}
% ======================================================================
    
    We work on a four-dimensional Lorentzian spacetime $(\mathcal{M},g_{\mu\nu})$. A detector is modeled as a timelike world-tube $\Sig \subset \mathcal{M}$, with intrinsic coordinates $\sigma^a$ ($a=1,2,3$), induced metric $h_{ab}$, and outward-pointing unit normal $n_\mu$.
    
    Let $J^\mu$ denote a conserved matter current,
    \begin{equation}
        \nabla_\mu J^\mu = 0 .
    \end{equation}
    
    \begin{definition}[Matter flux]
The matter flux through $\Sig$ is
\begin{equation}
    \F(\sigma) \equiv n_\mu J^\mu\big(X(\sigma)\big).
\end{equation}
    \end{definition}
    
    The total registered flux is
    \begin{equation}
        \mathcal{N} = \int_\Sig \dd^3\sigma\,\sqrt{|h|}\,\F(\sigma).
    \end{equation}

% ======================================================================
    \section{Event Currents and Discrete Clocks}\label{sec:eventcurrents}
% ======================================================================
    
    \begin{definition}[Effective event current]
An \emph{effective event current} is a vector field $\Jev$ satisfying
\begin{equation}
    \nabla_\mu \Jev \approx 0 ,
\end{equation}
where equality may be exact or hold up to controlled corrections on the scales of interest.
    \end{definition}
    
    \begin{remark}[Status of the event current]
The current $j^\mu_{\mathrm{ev}}$ is not assumed to be unique or fundamental. Rather, it represents a coarse-grained, detector-relevant counting current associated with a specified measurement context. Different admissible choices of $j^\mu_{\mathrm{ev}}$ correspond, in general, to different effective clock theories and therefore to different TOA statistics. This freedom is not hidden in the formalism: it is part of the model specification and is, in principle, experimentally distinguishable through the resulting clock correlators and arrival-time distributions.
    \end{remark}
    
    \begin{remark}[Concrete example]
In a relativistic QFT scattering experiment, $j^\mu_{\mathrm{ev}}$ may represent the effective flux of detector-registered interaction events, i.e.\ a spacetime current whose integral counts localized vertex/crossing events relevant to the measurement.
    \end{remark}
    
    \begin{definition}[Event count]
The associated event count through a hypersurface $\Sig$ is
\begin{equation}
    \Nev(\Sig) \equiv \int_\Sig \dd\Sigma_\mu\,\Jev .
\end{equation}
    \end{definition}
    
    \begin{proposition}
        If $\nabla_\mu \Jev = 0$, then $\Nev(\Sig)$ is invariant under smooth deformations of $\Sig$ that preserve its boundary.
    \end{proposition}
    
    \begin{proof}[Sketch]
The difference between the integrals over two such hypersurfaces is the flux of $\nabla_\mu j^\mu_{\mathrm{ev}}$ through the interpolating spacetime region. The claim then follows from Stokes' theorem and current conservation.
    \end{proof}
    
    The discrete-event clock is
    \begin{equation}
        T_{\mathrm{disc}} \equiv \Nev .
    \end{equation}
    
    \paragraph{Conceptual hierarchy.}
        The logical structure of the construction may be summarized schematically as
        \begin{center}
            \begin{tabular}{c c c}
                event current $j^\mu_{\mathrm{ev}}$
                & $\longrightarrow$ &
                discrete clock $T_{\mathrm{disc}} = N_{\mathrm{ev}}$ \\
                & &
                $\downarrow$ coarse-graining \\
                & &
                one-form $a^{(\ell)}_\mu=(j^{\mathrm{ev}}_\mu)_\ell$ \\
                & &
                $\downarrow$ integrability $\Omega^{\mathrm{ev}}_{\mu\nu}=0$ \\
                & &
                continuum clock field $T(x)$
            \end{tabular}
        \end{center}
        The event current is the ultraviolet object; event counts, clock one-forms, and continuum scalar clocks are progressively more restrictive effective descriptions.

% ======================================================================
    \section{Coarse-Graining, Integrability, and Clock Holonomy}
    \label{sec:clocklimit}
% ======================================================================
    
    Let $(\cdot)_\ell$ denote coarse-graining over a spacetime scale $\ell$, and define the coarse-grained event-current one-form
    \begin{equation}
        a^{(\ell)}_\mu(x) \equiv \big(j^{\mathrm{ev}}_\mu(x)\big)_\ell .
        \label{eq:clockoneform}
    \end{equation}
    Conservation and clock-field existence impose different conditions. Conservation,
    \begin{equation}
        \nabla_\mu j^\mu_{\mathrm{ev}}\approx 0,
    \end{equation}
    makes event counts stable under deformations of the counting surface. By contrast, the existence of a scalar continuum clock requires the one-form $a^{(\ell)}_\mu$ to be locally exact.
    
    We define the event-current vorticity, or clock-integrability obstruction, by
    \begin{equation}
        \Omega^{\mathrm{ev}}_{\mu\nu}
        \equiv
        \partial_\mu a^{(\ell)}_\nu
        -
        \partial_\nu a^{(\ell)}_\mu
        =
        2\partial_{[\mu}\big(j^{\mathrm{ev}}_{\nu]}\big)_\ell .
        \label{eq:eventvorticity}
    \end{equation}
    The integrable-clock sector is the sector in which
    \begin{equation}
        \Omega^{\mathrm{ev}}_{\mu\nu}=0
        \qquad\text{on the infrared patch of interest.}
        \label{eq:integrability}
    \end{equation}
    Under this condition one may define a local scalar clock field $T(x)$ by
    \begin{equation}
        \partial_\mu T(x) \equiv \frac{1}{\nu_0}\,a^{(\ell)}_\mu(x),
        \label{eq:clockgradient}
    \end{equation}
    where $\nu_0$ is an input conversion factor with dimensions of events per unit time.
    
    \begin{proposition}[Continuum clock criterion]
        On a simply connected infrared patch $U\subset\mathcal{M}$, a scalar clock field $T$ satisfying Eq.~\eqref{eq:clockgradient} exists locally if and only if
        \begin{equation}
            \Omega^{\mathrm{ev}}_{\mu\nu}=0
            \qquad \text{on } U .
        \end{equation}
    \end{proposition}
    
    \begin{proof}[Sketch]
        If $a^{(\ell)}_\mu=\nu_0\partial_\mu T$, then $\partial_\mu a^{(\ell)}_\nu-\partial_\nu a^{(\ell)}_\mu=0$ by equality of mixed partial derivatives. Conversely, if the one-form $a^{(\ell)}=a^{(\ell)}_\mu\dd x^\mu$ is closed on a simply connected patch, the Poincar\'e lemma implies that it is exact, $a^{(\ell)}=\nu_0\dd T$, for some scalar $T$ defined up to an additive constant.
    \end{proof}
    
    \begin{remark}[Conservation versus integrability]
        The distinction is essential. A divergence-free current need not be curl-free. A rotating or vortical congruence can conserve event number while failing to define a single-valued scalar time function. Thus Eq.~\eqref{eq:integrability} is not a mild technical afterthought; it is the primary condition for the continuum-clock limit.
    \end{remark}
    
    \subsection{Clock holonomy in the non-integrable sector}
    
    If $\Omega^{\mathrm{ev}}_{\mu\nu}\neq0$, there is in general no scalar clock field $T(x)$ whose gradient equals $a^{(\ell)}_\mu/\nu_0$. Nevertheless, the event current still defines a path-dependent clock functional. For a curve $\gamma$ from $x_i$ to $x_f$, define
    \begin{equation}
        \Delta T_\gamma
        \equiv
        \frac{1}{\nu_0}\int_\gamma a^{(\ell)}_\mu\,\dd x^\mu .
        \label{eq:pathclock}
    \end{equation}
    For two curves $\gamma_1$ and $\gamma_2$ with the same endpoints and bounding an oriented surface $S$, Stokes' theorem gives
    \begin{equation}
        \Delta T_{\gamma_1}-\Delta T_{\gamma_2}
        =
        \frac{1}{2\nu_0}
        \int_S
        \Omega^{\mathrm{ev}}_{\mu\nu}\,\dd S^{\mu\nu} .
        \label{eq:clockholonomy}
    \end{equation}
    We call this difference the \emph{clock holonomy}. It is the measurable obstruction to replacing the discrete/coarse-grained event-count clock by a single-valued continuum clock field.
    
    \begin{proposition}[Holonomy obstruction]
        A continuum clock is path-independent on a patch if and only if the clock holonomy \eqref{eq:clockholonomy} vanishes for all closed loops in that patch. Equivalently, $\Omega^{\mathrm{ev}}_{\mu\nu}=0$ locally and any remaining obstruction is global/topological.
    \end{proposition}
    
    \begin{remark}[Physical interpretation]
        Clock holonomy is not an additional ontology. It is the standard differential-geometric obstruction that appears when a one-form cannot be represented as the gradient of a scalar. In the present TOA setting it has a direct operational meaning: two admissible histories through the same detector region can accumulate different clock readings even when they share endpoints. This is precisely the kind of construction-dependent effect absent from a generic ``true time plus Gaussian noise'' model.
    \end{remark}
    
    \begin{proposition}[IR equivalence in the integrable sector]
        At scales $\gg \ell$, if $\Omega^{\mathrm{ev}}_{\mu\nu}=0$ and the detector region compares two nearby level sets of $T$, then event counts and clock differences match locally:
        \begin{equation}
            \Nev(\Sig)\approx \nu_0\int_\Sig \dd\Sigma_\mu\,\partial^\mu T
            =\nu_0\,\Delta T .
        \end{equation}
        For generic finite detectors in $3+1$ dimensions this equality is not a global identity; its domain of validity is the integrable local-clock regime specified above. Outside that regime, the path-dependent expression \eqref{eq:pathclock} and its holonomy \eqref{eq:clockholonomy} replace the scalar-clock description.
    \end{proposition}

% ======================================================================
    \section{Clock-Sector Effective Field Theory}\label{sec:clockeft}
% ======================================================================
    
    In the integrable sector of Section~\ref{sec:clocklimit}, one may describe the continuum clock by a local EFT. For the purpose of a local EFT expansion in a Minkowski background, we choose an inertial coordinate patch and expand the clock field around a monotonically increasing background solution,
    \begin{equation}
        T(x)=T_0(x)+\tau(x), \qquad T_0(x)=t,
    \end{equation}
    so that $\tau$ describes fluctuations around the reference clock profile.
    
    A minimal stable quadratic EFT is
    \begin{equation}
        S_\tau = \frac{1}{2}\int \dd^4x
        \left[
            (\partial_t\tau)^2
        - c_\tau^2(\nabla\tau)^2
        - \mu_\tau^2 \tau^2
        \right].
        \label{eq:clockEFT}
    \end{equation}
    The parameter $\mu_\tau$ sets a correlation scale in space and time.
    
    In Appendix~\ref{app:prop}, $\mu_\tau$ is shown to control the exponential suppression envelope of spacelike clock correlations. The use of that envelope in early-arrival estimates is model-dependent and is stated only as a heuristic consequence of the clock correlator.
    
    \begin{remark}[Correlation scale]
The mass parameter $\mu_\tau$ sets the characteristic temporal and spatial correlation scale of the clock field. As shown explicitly in Appendix~\ref{app:prop}, the same parameter controls the exponential suppression envelope of spacelike clock correlations. Any inference from that correlator envelope to an early-arrival tail is model-dependent and requires a specified detector--clock coupling.
    \end{remark}

% ======================================================================
    \section{Relational Time-of-Arrival Observable}\label{sec:toaobservable}
% ======================================================================
    
    \begin{definition}[Relational TOA distribution]
The TOA probability density with respect to the clock field $T$ is
\begin{equation}
    p(\Theta) \equiv
    \frac{1}{\mathcal{N}}
    \int_\Sig \dd^3\sigma\,\sqrt{|h|}\,
    \F(\sigma)\,
    \delta\!\big(T_\Sig(\sigma)-\Theta\big),
    \label{eq:TOA}
\end{equation}
where $T_\Sig(\sigma)\equiv T(X(\sigma))$.
    \end{definition}
    
    \paragraph{Why this is nontrivial.}
        Equation~\eqref{eq:TOA} is not a relabeling of an external time coordinate. It applies in the integrable sector where the coarse-grained event-current one-form admits a scalar clock field. The nontrivial construction-dependent content is the criterion for when such a scalar clock exists, Eq.~\eqref{eq:eventvorticity}, and the holonomy obstruction, Eq.~\eqref{eq:clockholonomy}, that replaces it in the non-integrable sector. When a clock EFT is specified, its correlators enter TOA statistics through the averaging in Eq.~\eqref{eq:pThetaAvg} and can produce broadened arrival-time distributions.
        
        This definition is covariant and makes detector dependence explicit through the geometry of $\Sig$, the effective flux functional $\F$, the admissible clock current, and the clock-sector correlators.

    \subsection{Path-conditioned TOA in the non-integrable sector}

        When $\Omega^{\mathrm{ev}}_{\mu\nu}\neq0$, the scalar clock field $T(x)$ used in Eq.~\eqref{eq:TOA} does not exist on the patch under consideration. The observable can nevertheless be defined once one specifies how the clock readout is transported from a reference event to the detector world-tube. Let $x_0$ be a reference event and let $\mathcal{P}=\{\gamma_\sigma\}_{\sigma\in\Sigma}$ denote a smooth family of curves, with each $\gamma_\sigma$ running from $x_0$ to $X(\sigma)\in\Sigma$. Define the path-conditioned clock reading
        \begin{equation}
            T_{\mathcal{P}}(\sigma)
            \equiv
            T_0+
            \frac{1}{\nu_0}
            \int_{\gamma_\sigma}
            a^{(\ell)}_\mu\,\dd x^\mu .
            \label{eq:pathClockOnDetector}
        \end{equation}
        The corresponding non-integrable TOA distribution is
        \begin{equation}
            p_{\mathcal{P}}(\Theta)
            \equiv
            \frac{1}{\mathcal{N}}
            \int_\Sigma \dd^3\sigma\,\sqrt{|h|}\,
            \F(\sigma)\,
            \delta\!\left(T_{\mathcal{P}}(\sigma)-\Theta\right).
            \label{eq:pathTOA}
        \end{equation}
        This expression reduces to Eq.~\eqref{eq:TOA} whenever $a^{(\ell)}_\mu=\nu_0\partial_\mu T$, because then $T_{\mathcal{P}}(\sigma)$ is independent of the path family up to an additive calibration constant.

        For two admissible path families $\mathcal{P}$ and $\mathcal{P}'$ with the same endpoints on $\Sigma$, the pointwise clock-readout shift is
        \begin{equation}
            T_{\mathcal{P}}(\sigma)-T_{\mathcal{P}'}(\sigma)
            =
            \frac{1}{2\nu_0}
            \int_{S_\sigma}
            \Omega^{\mathrm{ev}}_{\mu\nu}\,\dd S^{\mu\nu},
            \label{eq:pathFamilyShift}
        \end{equation}
        where $S_\sigma$ is any oriented surface bounded by the closed curve $\gamma_\sigma\circ\gamma_\sigma'^{-1}$. Thus the non-integrable sector yields a measurable effect only after a path-family specification, and the difference between such specifications is governed by the same holonomy obstruction defined in Eq.~\eqref{eq:clockholonomy}.

        \begin{remark}[Operational status of the path family]
            The path family $\mathcal{P}$ is not an extra hidden variable. It represents the operational protocol by which the clock sector transports or compares readouts across the detector region. In an integrable clock sector this protocol drops out; in a non-integrable sector it becomes part of the detector model and can in principle be constrained by comparing TOA statistics under different clock-routing or synchronization procedures.
        \end{remark}

        \paragraph{Schematic observability estimate.}
            The holonomy shift becomes operational when two clock-routing protocols enclose a nonzero effective area in a region with non-vanishing event-current vorticity. If $\Omega^{\mathrm{ev}}_{\mu\nu}$ is approximately constant over a characteristic oriented surface $S$ with effective area $A_{\mathrm{loop}}$, Eq.~\eqref{eq:pathFamilyShift} gives the parametric estimate
            \begin{equation}
                \left|\Delta T_{\mathcal{P}\mathcal{P}'}\right|
                \sim
                \frac{|\Omega^{\mathrm{ev}}|\,A_{\mathrm{loop}}}{2\nu_0}.
                \label{eq:holonomyEstimate}
            \end{equation}
            For two fixed, repeatable synchronization protocols $\mathcal{P}$ and $\mathcal{P}'$, the leading observable is a flux-weighted displacement of the mean TOA,
            \begin{equation}
                \Delta\langle\Theta\rangle_{\mathcal{P}\mathcal{P}'}
                =
                \frac{1}{\mathcal{N}}
                \int_{\Sigma}\dd^3\sigma\,\sqrt{|h|}\,
                \mathcal{F}(\sigma)
                \big[T_{\mathcal{P}}(\sigma)-T_{\mathcal{P}'}(\sigma)\big].
                \label{eq:holonomyMeanShiftGeneral}
            \end{equation}
            In the leading uniform-holonomy regime, where $T_{\mathcal{P}}(\sigma)-T_{\mathcal{P}'}(\sigma)$ is approximately constant over the support of $\mathcal{F}$ on $\Sigma$, this reduces to
            \begin{equation}
                \Delta\langle\Theta\rangle_{\mathcal{P}\mathcal{P}'}
                \simeq
                \Delta T_{\mathcal{P}\mathcal{P}'}.
                \label{eq:holonomyMeanShift}
            \end{equation}
            This deterministic protocol comparison contributes primarily to the mean, not to the variance.
            A variance contribution arises only when the path family is itself sampled from a specified ensemble $\mathcal{E}$ of synchronization protocols. In that case we define
            \begin{equation}
                \sigma_{\mathrm{hol}}^2
                \equiv
                \frac{1}{\mathcal{N}}
                \int_{\Sigma}\dd^3\sigma\,\sqrt{|h|}\,
                \mathcal{F}(\sigma)
                \operatorname{Var}_{\mathcal{P}\in\mathcal{E}}
                \!\left[T_{\mathcal{P}}(\sigma)\right],
                \label{eq:sigmaHolDefinition}
            \end{equation}
            If the pointwise root-mean-square loop area $A_{\mathrm{rms}}(\sigma)$ and the magnitude of $\Omega^{\mathrm{ev}}_{\mu\nu}$ vary slowly across the support of $\mathcal{F}$ on $\Sigma$, then the pointwise variance approximately factorizes out of the flux integral. In this leading uniform-holonomy regime, Eq.~\eqref{eq:sigmaHolDefinition} has the parametric scale
            \begin{equation}
                \sigma_{\mathrm{hol}}^2
                \sim
                \left(\frac{|\Omega^{\mathrm{ev}}|\,A_{\mathrm{rms}}}{2\nu_0}\right)^2,
                \label{eq:sigmaHolEstimate}
            \end{equation}
            where $A_{\mathrm{rms}}$ is the approximately flux-independent root-mean-square effective area enclosed by the protocol fluctuations in the ensemble.
            If the wavepacket-width fluctuations, the local clock-field fluctuations, and the clock-routing fluctuations are mutually uncorrelated, then the leading variance is
            \begin{equation}
                \operatorname{Var}(\Theta)
                \simeq
                \sigma_t^2+\sigma_{\tau,\ell}^2+\sigma_{\mathrm{hol}}^2.
                \label{eq:holonomyVariance}
            \end{equation}
            This independence assumption is appropriate, for example, when the routing protocol is fixed by a classical synchronization procedure independent of both the incoming quantum matter state and the microscopic clock-field fluctuation $\tau$. Thus the non-integrable sector is not merely diagnostic: once the clock-routing protocol, or an ensemble of such protocols, and the effective loop area are specified, the holonomy gives a direct parametric target for a mean TOA shift or an excess protocol-dependent variance.

% ======================================================================
    \section{Semiclassical Limit}
% ======================================================================
    
    In the semiclassical regime of sharply peaked wavepackets and weak clock fluctuations,
    \begin{equation}
        p(\Theta)\to \delta(\Theta-\Theta_{\mathrm{cl}}),
    \end{equation}
    where $\Theta_{\mathrm{cl}}$ matches the classical time-of-flight.

% ======================================================================
    \section{Worked Example in $1+1$ Dimensions}\label{sec:worked-example}
% ======================================================================
    
    We now give an explicit $1+1$D computation (wavepacket + clock correlator). This section is intended as a concrete longitudinal-sector template; the same construction can be extended to $3+1$D by restoring transverse detector integrals. Although the clock EFT in Eq.~\eqref{eq:clockEFT} was written in covariant $3+1$D form, the present worked example evaluates only its dimensionally reduced longitudinal sector, with transverse clock-mode contributions taken to be negligible for this illustrative computation. Accordingly, the momentum integrals and correlators below are $1+1$D quantities; the closed form \eqref{eq:explicitBroadening} should be read as the regulated variance of that longitudinal sector.
    
    \subsection{Setup: detector worldline and matter wavepacket}
        
        Work in Minkowski space with coordinates $(t,x)$ and metric $\eta=\mathrm{diag}(1,-1)$. Model the detector as the timelike worldline at fixed $x=x_D$:
        \begin{equation}
            \Sig:\quad X^\mu(\sigma)=(t=\sigma,\ x=x_D).
        \end{equation}
        For a $1+1$D detector ``world-tube'', the induced measure reduces to $\dd\sigma=\dd t$.
        
        Take a complex Klein--Gordon field $\psi$ of mass $m$ with conserved $U(1)$ current \cite{Peskin1995,Weinberg1995}:
        \begin{equation}
            J^\mu = \frac{i}{2m}\left(\psi^\ast \partial^\mu \psi - \psi \partial^\mu \psi^\ast\right),
            \qquad \partial_\mu J^\mu = 0.
            \label{eq:KGcurrent}
        \end{equation}
        In $1+1$D, the spatial flux at the detector is $J^x(t,x_D)$. We take the incoming direction to be $+x$ and define the registered flux by
        \begin{equation}
            \F(t)\equiv J^x(t,x_D)\,\mathbbm{1}_{\{J^x\ge 0\}}.
            \label{eq:positiveflux}
        \end{equation}
        This positive-flux restriction is an operational assumption, standard in TOA constructions \cite{Kijowski1974,Giannitrapani1997,Muga2008}, and is appropriate in the narrowband, predominantly right-moving regime considered below. It is not itself a fully covariant replacement for the signed world-tube flux; rather, it is a detector-frame reduction of the covariant flux functional in a regime with a preferred incoming direction. In that regime the negative-flux component is parametrically suppressed, so Eq.~\eqref{eq:positiveflux} modifies neither the normalization nor the convolution formula beyond subleading corrections. Outside that regime, backflow or bidirectional contributions must be treated explicitly through the full signed flux and detector geometry, and the simple positive-flux truncation need not be adequate.
        
        Let the positive-frequency wavepacket be
        \begin{equation}
            \psi(t,x)=\int_{-\infty}^{+\infty}\frac{\dd k}{\sqrt{2\pi}}\,
            a(k)\,e^{-i\omega_k t + i k x},
            \qquad \omega_k=\sqrt{k^2+m^2}.
            \label{eq:wavepacket}
        \end{equation}
        Choose $a(k)$ narrowly peaked at $k_0>0$ with width $\Delta k\ll k_0$. Then the packet propagates with group velocity
        \begin{equation}
            v_g=\frac{\partial \omega_k}{\partial k}\Big|_{k_0}=\frac{k_0}{\omega_{k_0}}\in(0,1).
        \end{equation}
    
    \subsection{Semiclassical approximation for the flux}
        
        Under the narrowband approximation, $J^x$ reduces to the familiar form
        \begin{equation}
            J^x(t,x)\approx v_g\,|\psi(t,x)|^2,
            \label{eq:Japprox}
        \end{equation}
        up to corrections of order $\cO((\Delta k/k_0)^2)$. A Gaussian $a(k)$ gives an approximately Gaussian arrival profile at $x=x_D$:
        \begin{equation}
            \F(t)\approx \F_0\,\exp\!\left[-\frac{(t-t_{\mathrm{cl}})^2}{2\sigma_t^2}\right],
            \qquad
            t_{\mathrm{cl}}=\frac{x_D-x_0}{v_g},
            \label{eq:fluxGaussian}
        \end{equation}
        where $x_0$ is the initial packet center and $\sigma_t\sim 1/(v_g\Delta k)$. Normalize $\mathcal{N}=\int \dd t\,\F(t)$.
    
    \subsection{Clock field on the detector and stationary correlator}
        
        Take the continuum clock field
        \begin{equation}
            T(t,x)=t+\tau(t,x).
        \end{equation}
        The relational TOA density \eqref{eq:TOA} becomes in $1+1$D
        \begin{equation}
            p(\Theta)=\frac{1}{\mathcal{N}}\int_{-\infty}^{+\infty}\dd t\,
            \F(t)\,
            \Big\langle \delta\!\big(t+\tau(t,x_D)-\Theta\big)\Big\rangle_{\tau},
            \label{eq:pThetaAvg}
        \end{equation}
        where $\langle\cdot\rangle_\tau$ denotes averaging over clock fluctuations.
        
        Assume the clock sector is in a stationary Gaussian state induced by the quadratic EFT \eqref{eq:clockEFT}. Then the single-point random variable $\tau_\ell(t,x_D)$ is Gaussian, with mean $0$ and variance
        \begin{equation}
            \sigma_{\tau,\ell}^2 \equiv \langle \tau_\ell(t,x_D)^2\rangle = C_\ell(0,0),
        \end{equation}
        where the two-point function is
        \begin{equation}
            C(\Delta t,\Delta x)\equiv \langle \tau(t,x)\tau(t+\Delta t,x+\Delta x)\rangle .
        \end{equation}
        Here $\tau_\ell$ denotes the clock field smeared over the detector's finite spacetime resolution scale $\ell$. This coarse-grained variance is the physically relevant quantity entering the TOA distribution. In a continuum free-field description, the unsmeared coincidence limit $C(0,0)$ is UV-sensitive in $1+1$ dimensions, so the theory should be understood as defined either with an explicit ultraviolet cutoff $\Lambda$ or, equivalently for present purposes, with a finite detector/coarse-graining scale $\ell$. Accordingly, the observable variance is really $\sigma_{\tau,\ell}^2=\sigma_{\tau,\ell}^2(\mu_\tau)$, or equivalently $\sigma_\tau^2(\Lambda,\mu_\tau)$ in cutoff language.
        
        In $1+1$D, for a gapped relativistic scalar ($c_\tau=1$) at equal position $\Delta x=0$, the detector-worldline correlator is timelike rather than spacelike. Its positive-frequency form is
        \begin{equation}
            C(\Delta t,0)
            =
            \int_{-\infty}^{+\infty}\frac{\dd k}{4\pi\omega_k}\,
            e^{-i\omega_k\Delta t},
            \qquad
            \omega_k=\sqrt{k^2+\mu_\tau^2} .
            \label{eq:timelikeCorrelator}
        \end{equation}
        For large $|\Delta t|$, stationary phase around $k=0$ gives an oscillatory algebraic envelope,
        \begin{equation}
            C(\Delta t,0)
            \sim
            \frac{\cos\!\big(\mu_\tau |\Delta t|+\varphi_0\big)}
            {\sqrt{\mu_\tau |\Delta t|}}
            \quad\text{for}\quad |\Delta t|\gg \mu_\tau^{-1},
            \label{eq:CtimelikeAsymptotic}
        \end{equation}
        up to convention-dependent phases and algebraic prefactors. Thus $\mu_\tau^{-1}$ sets the characteristic oscillation scale of timelike clock correlations, not an exponential damping length. Exponential suppression is instead a spacelike-correlation statement, derived explicitly in Appendix~\ref{app:prop}.
        
        To make the regulated variance explicit, let $K_\ell(t,x)$ be a normalized detector-resolution kernel and define
        \begin{equation}
            \tau_\ell(t,x)
            =
            \int \dd t'\dd x'\,K_\ell(t-t',x-x')\,\tau(t',x') .
            \label{eq:tausmear}
        \end{equation}
        Then
        \begin{equation}
            \sigma_{\tau,\ell}^{2}
            =
            \int \dd^2y\,\dd^2z\,
            K_\ell(y)K_\ell(z)\,G^{+}(y-z) .
            \label{eq:regulatedvariance}
        \end{equation}
        Equivalently, in a translation-invariant positive-frequency representation,
        \begin{equation}
            \sigma_{\tau,\ell}^{2}
            =
            \int\frac{\dd k}{2\pi}\,
            \frac{\left|\widetilde K_\ell(\omega_k,k)\right|^2}{2\omega_k},
            \qquad
            \omega_k=\sqrt{k^2+\mu_\tau^2} .
            \label{eq:regulatedvarianceMomentum}
        \end{equation}
        This equation is the quantity that must be supplied by a concrete clock model or detector calibration. For smooth kernels $K_\ell$ it is finite and depends on the two physical scales $\mu_\tau^{-1}$ and $\ell$; without such smearing the coincidence limit is not an observable.

        As a concrete evaluation, choose a Gaussian detector-resolution kernel whose on-shell Fourier weight is
        \begin{equation}
            \bigl|\widetilde K_\ell(\omega_k,k)\bigr|^2
            =
            \exp\!\big[-\ell^2(\omega_k^2+k^2)\big],
            \qquad
            \omega_k^2=k^2+\mu_\tau^2 .
            \label{eq:gaussianKernelWeight}
        \end{equation}
        Here $\widetilde K_\ell(\omega_k,k)$ denotes the positive-frequency spectral weight of the smearing kernel evaluated on the mass shell $\omega=\omega_k$, not an independent off-shell Fourier transform. The full off-shell transform $\widetilde K_\ell(\omega,k)$ is not used as a separate dynamical object below; only its on-shell restriction weights the positive-frequency spectral density. A normalized Gaussian position-space resolution kernel $K_\ell(t,x)\propto\exp[-(t^2+x^2)/(2\ell^2)]$, restricted to the positive-frequency clock spectrum, gives precisely this on-shell weight up to normalization conventions absorbed into the definition of $K_\ell$.
        Then Eq.~\eqref{eq:regulatedvarianceMomentum} gives
        \begin{equation}
            \sigma_{\tau,\ell}^{2}
            =
            \frac{e^{-\mu_\tau^2\ell^2}}{4\pi}
            \int_{-\infty}^{+\infty}
            \frac{e^{-2\ell^2 k^2}}{\sqrt{k^2+\mu_\tau^2}}\,\dd k .
            \label{eq:gaussianVarianceIntegral}
        \end{equation}
        Using $k=\mu_\tau\sinh u$ and
        $\int_0^\infty e^{-z\cosh t}\dd t=K_0(z)$, the integral evaluates to
        \begin{equation}
            \boxed{
            \sigma_{\tau,\ell}^{2}
            =
            \frac{1}{4\pi}
            K_0(\mu_\tau^2\ell^2)
            }
            \label{eq:gaussianVarianceClosed}
        \end{equation}
        for this smearing convention. Therefore the explicitly regulated broadening in the Gaussian clock model is
        \begin{equation}
            \mathrm{Var}(\Theta)
            =
            \sigma_t^2+\frac{1}{4\pi}K_0(\mu_\tau^2\ell^2),
            \label{eq:explicitBroadening}
        \end{equation}
        under the same spectator-clock and Gaussian-flux assumptions used below.

        The limiting behavior is transparent. For $\mu_\tau\ell\ll1$,
        \begin{equation}
            \sigma_{\tau,\ell}^{2}
            \sim
            \frac{1}{4\pi}
            \left[
            -\ln\left(\frac{\mu_\tau^2\ell^2}{2}\right)-\gamma_E
            \right],
            \label{eq:smallSmearingLimit}
        \end{equation}
        while for $\mu_\tau\ell\gg1$,
        \begin{equation}
            \sigma_{\tau,\ell}^{2}
            \sim
            \frac{1}{4\pi}
            \sqrt{\frac{\pi}{2\mu_\tau^2\ell^2}}
            e^{-\mu_\tau^2\ell^2}.
            \label{eq:largeSmearingLimit}
        \end{equation}
        Thus the clock broadening is logarithmically UV-sensitive as the detector resolution is removed, but becomes exponentially suppressed once the smearing scale exceeds the clock correlation length.
    
    \subsection{Closed-form TOA distribution: convolution formula}
        
        For a stationary Gaussian variable $X=\tau_\ell(t,x_D)$ with variance $\sigma_{\tau,\ell}^2$,
        \begin{equation}
            \left\langle \delta(t+X-\Theta)\right\rangle_X
            =\frac{1}{\sqrt{2\pi\sigma_{\tau,\ell}^2}}
            \exp\!\left[-\frac{(\Theta-t)^2}{2\sigma_{\tau,\ell}^2}\right].
        \end{equation}
        To obtain a closed convolution form, we further assume that, at leading order, the local clock fluctuation $\tau(t,x_D)$ is statistically independent of the matter-flux profile $\mathcal{F}(t)$. Physically, this corresponds to a weak-coupling or spectator-clock regime in which the clock sector broadens the readout without appreciably backreacting on the matter current. If clock--matter correlations are appreciable, mixed correlators contribute and the simple convolution formula below is replaced by a more general conditional-average expression.
        
        Under the Gaussian and leading-order independence assumptions just stated, Eq.~\eqref{eq:pThetaAvg} becomes the convolution
        \begin{equation}
            p(\Theta)=\frac{1}{\mathcal{N}}
            \int_{-\infty}^{+\infty}\dd t\,
            \F(t)\,
            \frac{1}{\sqrt{2\pi\sigma_{\tau,\ell}^2}}
            \exp\!\left[-\frac{(\Theta-t)^2}{2\sigma_{\tau,\ell}^2}\right].
            \label{eq:convolution}
        \end{equation}
        
        If $\F(t)$ is Gaussian as in \eqref{eq:fluxGaussian}, the convolution is again Gaussian:
        \begin{equation}
            p(\Theta)=\frac{1}{\sqrt{2\pi(\sigma_t^2+\sigma_{\tau,\ell}^2)}}
            \exp\!\left[-\frac{(\Theta-t_{\mathrm{cl}})^2}{2(\sigma_t^2+\sigma_{\tau,\ell}^2)}\right].
            \label{eq:pGaussian}
        \end{equation}
        Thus the clock sector produces a \emph{predictive broadening}:
        \begin{equation}
            \mathrm{Var}(\Theta)=\sigma_t^2+\sigma_{\tau,\ell}^2,
            \qquad
            \sigma_{\tau,\ell}^2=C_\ell(0,0).
            \label{eq:varianceAddition}
        \end{equation}
        For the Gaussian smearing choice in Eq.~\eqref{eq:gaussianKernelWeight}, this becomes the explicit expression \eqref{eq:explicitBroadening}.

        \paragraph{Falsifiability.}
            Given an experimentally or numerically specified flux profile $\mathcal{F}(t)$ and a clock variance $\sigma_{\tau,\ell}^2$ determined independently from the clock sector---for example through Eq.~\eqref{eq:regulatedvarianceMomentum} once $\mu_\tau$ and $\ell$ are fixed by detector calibration or separate clock-correlation measurements---the convolution formula~\eqref{eq:convolution} predicts a unique TOA distribution without additional fitting parameters. Deviations from the predicted variance $\mathrm{Var}(\Theta)=\sigma_t^2+\sigma_{\tau,\ell}^2$ or from the expected crossover behavior at $\Delta t\sim\mu_\tau^{-1}$ would directly falsify the assumed clock sector or its effective field description.

            \begin{remark}[Gaussian tails vs exponential envelopes]
The Gaussian form of $p(\Theta)$ arises from the \emph{local, one-point} statistics of the clock readout at the detector. Exponential behavior enters instead through the \emph{spacetime propagation} of clock correlations: early-arrival contributions requiring spacelike correlations are suggested, in models with the appropriate detector--clock coupling, to inherit an exponential \emph{envelope} set by $\mu_\tau^{-1}$ from the spacelike correlator derived in Appendix~\ref{app:prop}. These two effects are conceptually distinct and govern different regimes of the TOA distribution.
            \end{remark}

    \subsection{Discrete-event ultraviolet corrections}
    
    If the UV clock is an event-count $T_{\mathrm{disc}}=\Nev$, then at time resolution $\Delta t \lesssim \ell$ one expects renewal-type corrections:
    \begin{equation}
        p(\Theta)\;=\;p_{\mathrm{IR}}(\Theta)\;+\;\delta p_{\mathrm{UV}}(\Theta),
    \end{equation}
    where $p_{\mathrm{IR}}$ is the continuum prediction (e.g. \eqref{eq:pGaussian}) and $\delta p_{\mathrm{UV}}$ carries integer-step features and non-Gaussianity that vanish upon coarse-graining $\Delta t\gg \ell$.
    
    \medskip
    \noindent
    \textit{Bridge to a concrete realization.}
    So far the discussion has been model-independent: $J^\mu$ and $j^\mu_{\mathrm{ev}}$ are treated as effective conserved structures. The next section simply records one explicit realization in which $j^\mu_{\mathrm{ev}}$ is generated by defect worldlines; none of the preceding definitions rely on that choice.

% ======================================================================
    \section{Example Microscopic Realization: Defect Worldlines}
    \label{sec:defectrealization}
% ======================================================================
    
    The relational TOA construction above is model-independent: it requires only (i) a conserved matter flux $J^\mu$ through the detector world-tube $\Sigma$ and (ii) a conserved (or approximately conserved) event current $j^\mu_{\mathrm{ev}}$ whose integral defines the ultraviolet event-count clock. Here we record one concrete realization in an effective-medium/defect setting, presented as an illustration that the general framework accommodates explicit microscopic input; none of the preceding definitions or results depend on this particular choice.
    
    \subsection{Defect worldline current}
        
        Consider $N$ stable, localized defects with center worldlines $z_k^\mu(\tau)$. Define the distributional defect current
        \begin{equation}
            j^\mu_{\mathrm{ev}}(x)
            \;\equiv\;
            \sum_{k=1}^{N} q_k \int \dd\tau\,
            \frac{\dd z_k^\mu}{\dd\tau}\,
            \delta^{(4)}\!\big(x-z_k(\tau)\big),
            \label{eq:defectCurrent}
        \end{equation}
        where $q_k\in\mathbb{Z}$ is an integer-valued defect charge. In the absence of creation/annihilation events in the bulk, this current is conserved in the distributional sense,
        \begin{equation}
            \nabla_\mu j^\mu_{\mathrm{ev}}(x)=0,
        \end{equation}
        and the associated event count through a detector world-tube $\Sigma$ is
        \begin{equation}
            N_{\mathrm{ev}}(\Sigma)\equiv\int_{\Sigma}\dd\Sigma_\mu\, j^\mu_{\mathrm{ev}}.
        \end{equation}
        Operationally, $N_{\mathrm{ev}}(\Sigma)$ counts (with multiplicity $q_k$) the number of defect worldlines intersecting the detector region.
        
        In a coarse-grained ensemble with effective defect-line density $n_{\mathrm d}$, mean charge $q$, and characteristic crossing speed $v_{\mathrm d}$ in the detector frame, the conversion factor $\nu_0$ may be calibrated parametrically as the mean event-crossing rate,
        \begin{equation}
            \nu_0 \sim n_{\mathrm d}\,q\,v_{\mathrm d}\,A_{\mathrm D},
            \label{eq:nu0DefectCalibration}
        \end{equation}
        for a detector cross-section $A_{\mathrm D}$, up to geometry and acceptance factors. More generally, $\nu_0$ is fixed by an independent clock-rate calibration: it converts event counts into the continuum time unit used by the infrared clock field.
    
    \subsection{Continuum clock field from coarse-graining}
        
        At scales large compared to a coarse-graining length $\ell$, the distributional current \eqref{eq:defectCurrent} may be smoothed into a one-form $a^{(\ell)}_\mu=(j^{\mathrm{ev}}_\mu)_\ell$. A smooth scalar clock field exists only if the integrability condition of Section~\ref{sec:clocklimit} holds:
        \begin{equation}
            \Omega^{\mathrm{ev}}_{\mu\nu}
            =
            \partial_\mu a^{(\ell)}_\nu-\partial_\nu a^{(\ell)}_\mu
            \approx 0 .
            \label{eq:defectIntegrability}
        \end{equation}
        When this condition is satisfied, define
        \begin{equation}
            \partial_\mu T(x)
            \;\equiv\;
            \frac{1}{\nu_0}\,a^{(\ell)}_\mu(x),
            \label{eq:TFromEvents}
        \end{equation}
        where $\nu_0$ is a conversion factor. This yields the continuum-clock (IR) description of the event-count (UV) clock. When Eq.~\eqref{eq:defectIntegrability} fails, the same defect current defines the path-dependent clock functional \eqref{eq:pathclock}, the holonomy \eqref{eq:clockholonomy}, and the path-conditioned TOA distribution \eqref{eq:pathTOA}, rather than a single-valued scalar clock.

        A simple geometric criterion follows if the smoothed defect ensemble is represented by
        \begin{equation}
            a^{(\ell)}_\mu \simeq n_{\mathrm d}\,q\,u_\mu,
            \label{eq:defectLaminarAnsatz}
        \end{equation}
        where $u_\mu$ is the effective timelike velocity one-form of the defect congruence. Then
        \begin{equation}
            \Omega^{\mathrm{ev}}_{\mu\nu}
            =
            \partial_\mu(n_{\mathrm d}q u_\nu)
            -
            \partial_\nu(n_{\mathrm d}q u_\mu).
            \label{eq:defectOmega}
        \end{equation}
        If $n_{\mathrm d}q$ varies slowly on the coarse-graining scale, integrability reduces approximately to
        \begin{equation}
            \partial_{[\mu}u_{\nu]}\approx0,
            \label{eq:defectHypersurfaceOrthogonal}
        \end{equation}
        i.e. the defect worldlines must form an approximately laminar, hypersurface-orthogonal congruence. A randomly oriented or rotating ensemble of worldlines is generically non-integrable and therefore belongs to the holonomy sector rather than the scalar-clock sector.
    
    \subsection{Foliation interpretation}
        
        If the macroscopic medium admits a hypersurface-orthogonal timelike congruence $u_\mu$ (a foliation), and if $\nabla_\mu T$ is everywhere timelike in the region of interest, then one may parametrize it by a scalar potential $T$ via
        \begin{equation}
            u_\mu \;=\;\frac{\nabla_\mu T}{\sqrt{g^{\alpha\beta}\nabla_\alpha T\,\nabla_\beta T}}.
            \label{eq:uFromT}
        \end{equation}
        In this interpretation, the TOA observable $p(\Theta)$ defined in Eq.~\eqref{eq:TOA} becomes the distribution of detector flux binned by the physical clock reading $T$ on $\Sigma$. The foliation is therefore not merely interpretational: it is the geometric representation of the integrable-clock sector.
    
    \subsection{Interpretational note (optional)}
        
        The localized defects may be viewed as topological excitations of an effective medium, with the integer charges $q_k$ playing the role of quantized winding numbers or topological defect charges. This is standard language in the theory of topological defects in ordered media \cite{Peskin1995,Weinberg1995}.

% ======================================================================
    \section{Discussion and Outlook}\label{sec:discussion}
% ======================================================================
    
    \paragraph{What the event current represents (and what it does not).}
        In the main text, the event current $j^\mu_{\mathrm{ev}}$ is introduced as an (approximately) conserved quantity whose integral counts \emph{localized, robust events} relevant to a detection process. At this level of generality it is deliberately agnostic: the ``events'' might be particle crossings of a detector region, interaction vertices effectively registered by an apparatus, defect/excitation worldlines intersecting $\Sigma$, or other countable processes whose total number is stable under smooth deformations of $\Sigma$ that keep the boundary fixed.
        
        The minimal count-level assumption is that, on the scales of interest, the relevant events can be summarized by a current with $\nabla_\mu j^\mu_{\mathrm{ev}}\approx 0$ (exactly or with controlled violations). A stronger assumption is needed only when one wants a scalar continuum clock: the coarse-grained one-form must also be integrable, $\Omega^{\mathrm{ev}}_{\mu\nu}=0$. Section~\ref{sec:defectrealization} provides one concrete realization in terms of defect worldlines in an effective medium, but the relational TOA definition itself does not depend on adopting that ontology. Put differently, the construction is compatible with many microscopic stories; it only requires that they reduce to an admissible effective event-count description on appropriate scales.
        
        \medskip
        \noindent
        \textbf{Predictive content and where it enters.}
        Once a clock sector is specified, its correlators $C(\Delta t,\Delta x)$ become physical input rather than interpretational decoration. In the simplest (Gaussian) setting they determine the amount of arrival-time broadening through local clock variance, while massive clock correlations motivate a model-dependent exponential envelope for contributions that would require correlations across a spacelike separation (Appendix~\ref{app:prop}). In the non-integrable sector, different clock-path families can shift the readout according to the holonomy formula \eqref{eq:pathFamilyShift}, with the scale estimated in Eq.~\eqref{eq:holonomyEstimate}. At sufficiently fine time resolution, the ultraviolet event-count description can additionally produce discrete/renewal features that are washed out in the infrared after coarse-graining.

        \medskip
        \noindent
        \textbf{Admissibility of event currents.}
        The present formalism does not claim that every approximately conserved current is equally natural in a given experiment. The selection of an admissible $j^\mu_{\mathrm{ev}}$ is part of the effective model-building step and should be constrained by the detector context: locality, robustness under coarse-graining, compatibility with the detector's counting statistics, and empirical stability across changes of resolution. In this sense, the framework is not falsified or confirmed by a current in the abstract, but by whether a specified admissible current yields a clock sector whose correlators and TOA statistics agree with observation. A systematic classification of admissible event currents is left for future work.
    
    \subsection*{Idealized measurement scenarios}
    
    A minimal scenario is a relativistic wavepacket experiment in which the detector couples (directly or effectively) to a clock degree of freedom with a finite correlation time. The observable consequence is that the TOA statistics need not coincide with those obtained from the matter wavepacket alone. In the Gaussian clock limit discussed in Section~\ref{sec:worked-example}, the predicted broadening
    \begin{equation}
        \mathrm{Var}(\Theta)=\sigma_t^2+\sigma_{\tau,\ell}^2
    \end{equation}
    would appear as an excess arrival-time spread beyond the semiclassical flux width $\sigma_t$.
    
    From a practical standpoint, the framework is well-suited to numerical study. Given a flux profile $\mathcal{F}(t)$ and a model (or measurement-informed ansatz) for the clock correlator $C(\Delta t,\Delta x)$, the TOA density follows directly from the convolution formula, Eq.~\eqref{eq:convolution}. This makes the clock-sector parameters (e.g.\ $c_\tau$ and $\mu_\tau$ in the quadratic EFT) quantitatively testable in toy models and in lattice or continuum simulations.

% ======================================================================
    \section{Conclusion}
% ======================================================================
    
    We defined time-of-arrival (TOA) as a \emph{relational}, manifestly covariant observable: detector crossings are binned neither by an externally supplied time coordinate nor by postulating a universal time operator, but by a \emph{physical clock} evaluated on the detector world-tube $\Sigma$. The construction uses a conserved matter flux and an event-count clock sector, but it also makes explicit that conservation alone is insufficient for continuum time: a scalar clock exists only in the integrable sector where the coarse-grained event-current one-form is locally exact. The construction uses two ingredients:
    
    \smallskip
    \noindent
    (i) a conserved matter current $J^\mu$ whose contraction with the unit normal defines the detector flux $\mathcal{F}=n_\mu J^\mu$, and
    
    \noindent
    (ii) an event current $j^\mu_{\mathrm{ev}}$ whose integral defines a discrete event-count clock. The choice of $j^\mu_{\mathrm{ev}}$ is part of the effective model specification rather than an unobservable convention: different admissible event currents define different clock sectors and are, in principle, empirically distinguishable through the resulting TOA statistics.
    
    \smallskip
    A controlled coarse-graining over a scale $\ell$ then yields an infrared clock field $T(x)$ only when
    \begin{equation}
        \Omega^{\mathrm{ev}}_{\mu\nu}
        =
        \partial_\mu\big(j^{\mathrm{ev}}_\nu\big)_\ell
        -
        \partial_\nu\big(j^{\mathrm{ev}}_\mu\big)_\ell
        =0 .
    \end{equation}
    In that sector,
    \begin{equation}
        \partial_\mu T(x)=\frac{1}{\nu_0}\big(j_{\mathrm{ev},\mu}(x)\big)_\ell,
    \end{equation}
    so that discrete counts and continuum time arise as UV and IR descriptions of the same underlying structure. When $\Omega^{\mathrm{ev}}_{\mu\nu}\neq0$, the clock is path-dependent and its holonomy is
    \begin{equation}
        \Delta T_{\gamma_1}-\Delta T_{\gamma_2}
        =
        \frac{1}{2\nu_0}
        \int_S
        \Omega^{\mathrm{ev}}_{\mu\nu}\,\dd S^{\mu\nu} .
    \end{equation}
    
    In the integrable sector, the resulting scalar-clock TOA density is
    \begin{equation}
        p(\Theta)=
        \frac{1}{\mathcal{N}}
        \int_{\Sigma}\dd^3\sigma\,\sqrt{|h|}\,
        \mathcal{F}(\sigma)\,
        \delta\!\big(T_\Sigma(\sigma)-\Theta\big).
    \end{equation}
    In the non-integrable sector this is replaced by the path-conditioned density
    \begin{equation}
        p_{\mathcal{P}}(\Theta)=
        \frac{1}{\mathcal{N}}
        \int_{\Sigma}\dd^3\sigma\,\sqrt{|h|}\,
        \mathcal{F}(\sigma)\,
        \delta\!\big(T_{\mathcal{P}}(\sigma)-\Theta\big),
    \end{equation}
    where differences between admissible path families are controlled by the holonomy integral. In both cases the observable depends on detector microphysics only through explicit structures: the geometry of $\Sigma$, the effective flux $J^\mu$, the admissible event current, the clock-path family when needed, and the clock-sector correlators. Because the construction never requires a self-adjoint operator conjugate to the Hamiltonian, it sidesteps Pauli-type obstructions at the outset. In the semiclassical regime it reproduces the usual time-of-flight notion.

    In the worked $1+1$D example, a narrowband incoming wavepacket together with a stationary Gaussian clock sector reduces $p(\Theta)$ to a convolution of the semiclassical flux profile with the local clock readout distribution, producing a concrete prediction in the integrable spectator-clock regime:
    \begin{equation}
        \mathrm{Var}(\Theta)=\sigma_t^2+\sigma_{\tau,\ell}^2.
    \end{equation}
    In the non-integrable sector, an independently sampled ensemble of clock-path families adds the holonomy contribution defined in Eq.~\eqref{eq:sigmaHolDefinition}, giving the leading independent-fluctuation estimate
    \begin{equation}
        \mathrm{Var}(\Theta)\simeq
        \sigma_t^2+\sigma_{\tau,\ell}^2+\sigma_{\mathrm{hol}}^2 .
    \end{equation}
    For an explicit Gaussian detector-resolution kernel, the regulated clock variance evaluates to
    \begin{equation}
        \sigma_{\tau,\ell}^2=
        \frac{1}{4\pi}K_0(\mu_\tau^2\ell^2),
    \end{equation}
    with logarithmic sensitivity as $\mu_\tau\ell\to0$ and exponential suppression when $\mu_\tau\ell\gg1$. Beyond this local effect, Appendix~\ref{app:prop} shows that massive clock correlators are exponentially suppressed at spacelike separations with envelope scale $\mu_\tau^{-1}$ (up to algebraic prefactors), motivating model-dependent early-arrival envelope estimates without implying superluminal signaling.
    
    Relative to mainstream TOA formalisms, the point here is not that detector models or POVMs fail operationally, but that the relational definition isolates covariant field-functionals whose dependence on measurement details is pushed into explicit, modelable structures: conserved currents, integrability conditions, clock holonomies, clock-path families in the non-integrable sector, and a clock sector with its own dynamics. This reframes detector dependence as a UV$\rightarrow$IR matching problem---event counts to continuum time when the event-current one-form is integrable, or event counts to path-conditioned clock readings when it is not.
    
    \subsection*{Next steps}
        Immediate extensions include: (i) a $3+1$D worked example for finite-size detectors to quantify geometric dependence; (ii) non-Gaussian clock dynamics beyond the quadratic EFT to predict higher-cumulant signatures in $p(\Theta)$; and (iii) a systematic treatment of renewal-type UV corrections at temporal resolution $\Delta t\lesssim \ell$ and their crossover to the continuum limit. In particular, a numerical study of Eq.~\eqref{eq:convolution} as $\mu_\tau$ is varied would directly probe the regulated broadening \eqref{eq:explicitBroadening}, while a parallel study of Eq.~\eqref{eq:pathTOA} would probe holonomy-induced path dependence in non-integrable clock sectors.

% ======================================================================
        \appendix
    \section*{Appendix}
% ======================================================================
    \section{Exponential Suppression from the Clock Propagator}
    \label{app:prop}
% ======================================================================
    
    This appendix derives the exponential suppression scale governing spacelike clock correlations and motivates the corresponding exponential envelope expected in early-arrival tails for models in which such tails are controlled by clock-sector correlations. The derivation uses standard massive-scalar correlators and Bessel function asymptotics \cite{Peskin1995,Weinberg1995,GradshteynRyzhik}.
    
    \subsection{Clock correlators from the quadratic EFT}
        
        Consider the quadratic clock action in flat $1+1$D with $c_\tau=1$:
        \begin{equation}
            S_\tau=\frac12\int \dd^2x\left[(\partial_t\tau)^2-(\partial_x\tau)^2-\mu_{\tau}^{\,2}\tau^2\right].
        \end{equation}
        The Wightman function $\Gplus(x)$ for a free massive scalar is a Lorentz-invariant function of $s^2=(\Delta t)^2-(\Delta x)^2$:
        \begin{equation}
            \Gplus(\Delta t,\Delta x)=\langle \tau(t,x)\tau(t+\Delta t,x+\Delta x)\rangle.
        \end{equation}
        
        For spacelike separations $(\Delta x)^2-(\Delta t)^2>0$, standard results give (see, e.g., representations in \cite{Weinberg1995} and integral tables):
        \begin{equation}
            \Gplus(\Delta t,\Delta x)
            =\frac{1}{2\pi}\Kzero\!\left(\mu_\tau\,\sqrt{(\Delta x)^2-(\Delta t)^2}\right)
            \quad\text{(spacelike)}.
            \label{eq:WightmanK0}
        \end{equation}
        While conventions vary by factors and $i0^+$ prescriptions, the key fact is that the spacelike decay is controlled by $\Kzero$.
    
    \subsection{Bessel asymptotics and exponential decay}
        
        For large argument $z\gg 1$, the modified Bessel function satisfies
        \begin{equation}
            \Kzero(z)\sim \sqrt{\frac{\pi}{2z}}\,e^{-z}\left[1+\cO\!\left(\frac{1}{z}\right)\right].
            \label{eq:K0asym}
        \end{equation}
        Therefore, for spacelike separations with invariant distance
        \begin{equation}
            \rho \equiv \sqrt{(\Delta x)^2-(\Delta t)^2}\gg \mu_\tau^{-1},
        \end{equation}
        Eq.~\eqref{eq:WightmanK0} yields the exponential suppression
        \begin{equation}
            \Gplus(\Delta t,\Delta x)\;\propto\;
            \frac{e^{-\mu_\tau \rho}}{\sqrt{\mu_\tau \rho}}\,
            \left[1+\cO\!\left((\mu_\tau\rho)^{-1}\right)\right].
            \label{eq:spacelikeExp}
        \end{equation}
        
        \begin{proposition}[Spacelike correlator envelope]
For the quadratic clock EFT, the spacelike two-point function is exponentially suppressed with characteristic envelope scale $\mu_\tau^{-1}$ as in \eqref{eq:spacelikeExp}.
        \end{proposition}
        
        \begin{remark}
            This mechanism is directly analogous to standard QFT statements: massive fields do not transmit strong correlations outside the light cone; the residual correlations decay as $e^{-\mu \rho}$ (up to power-law prefactors). This is a correlator statement and does not imply superluminal signaling.
        \end{remark}
    
    \subsection{Heuristic implication for early-arrival tails}
        
        We now state a physically motivated inference rather than a theorem. In any TOA model in which contributions to sufficiently early detector readout require clock-sector correlations across an effectively spacelike interval, the spacelike suppression \eqref{eq:spacelikeExp} suggests an exponential envelope of the form
        \begin{equation}
            p_{\mathrm{early}}(\Delta)\;\lesssim\;A(\Delta)\,\exp\!\big(-\mu_\tau\,\rho(\Delta)\big),
            \label{eq:earlyBound}
        \end{equation}
        where $\Delta$ parameterizes the degree of earliness and $\rho(\Delta)$ denotes the corresponding spacelike separation scale.
        
        Here $A(\Delta)$ is left unspecified; in concrete realizations it must be derived from the detailed coupling between the clock sector and the TOA statistic. Accordingly, Eq.~\eqref{eq:earlyBound} should be read as a heuristic envelope statement induced by the clock correlator, not as a model-independent theorem for all TOA constructions.
    
    \subsection{Retarded Green's function and locality}
        
        For completeness, the retarded propagator solves
        \begin{equation}
            (\Box+\mu_\tau^2)\,\Gret(x)=\delta^{(2)}(x),
            \qquad \Gret(x)=0\ \text{for}\ t<0.
        \end{equation}
        Local detector couplings to $\tau$ then inherit causal support from $\Gret$, while spacelike Wightman correlations remain exponentially suppressed as above. This distinguishes \emph{correlation tails} from \emph{signal propagation}, consistent with standard relativistic QFT \cite{Peskin1995,Weinberg1995}.

% ======================================================================
% Back matter (FoP / Springer)
% ======================================================================
        
        \backmatter
        
        \bmhead{Acknowledgements}
        Not applicable.
        
        \bmhead{Declarations}
    
    \subsection*{Funding}
        Not applicable.
    
    \subsection*{Competing interests}
        On behalf of all authors, the corresponding author states that there is no conflict of interest.
    
    \subsection*{Ethics approval}
        Not applicable.
    
    \subsection*{Consent to participate}
        Not applicable.
    
    \subsection*{Consent for publication}
        Not applicable.
    
    \subsection*{Data availability}
        No datasets were generated or analyzed during the current study.
    
    \subsection*{Code availability}
        Not applicable.
    
    \subsection*{Author contributions}
        O.I. conceived the study, developed the analysis, performed the derivations, and wrote the manuscript.

% =========================
% References
% =========================
        \begin{thebibliography}{99}
            
            \bibitem{Pauli1980}
            W.~Pauli,
            \textit{General Principles of Quantum Mechanics},
            Springer (1980).
            
            \bibitem{Kijowski1974}
            J.~Kijowski,
            On the time operator in quantum mechanics and the Heisenberg uncertainty relation,
            \textit{Rep.\ Math.\ Phys.} \textbf{6}, 361 (1974).
            DOI: 10.1016/0034-4877(74)90051-5.
            
            \bibitem{Giannitrapani1997}
            G.~Giannitrapani,
            Positive-operator-valued measures for arrival times and time operators,
            \textit{Int.\ J.\ Theor.\ Phys.} \textbf{36}, 1575 (1997).
            DOI: 10.1007/BF02435766.
            
            \bibitem{Muga2008}
            J.~G.~Muga, R.~Sala Mayato, and I.~L.~Egusquiza (eds.),
            \textit{Time in Quantum Mechanics, Vol.~1},
            2nd ed., Lecture Notes in Physics \textbf{734},
            Springer (2008).
            DOI: 10.1007/978-3-540-73473-4.
            
            \bibitem{ColosiOeckl2025TimeOperatorTOA}
            D.~Colosi and R.~Oeckl,
            A Time Operator and the Time-Of-Arrival Problem in Quantum Field Theory,
            \textit{Foundations of Physics} \textbf{55}, 56 (2025).
            DOI: 10.1007/s10701-025-00866-x.
            
            \bibitem{Peskin1995}
            M.~E.~Peskin and D.~V.~Schroeder (1995),
            \textit{An Introduction to Quantum Field Theory},
            Addison-Wesley.
            
            \bibitem{Weinberg1995}
            S.~Weinberg (1995),
            \textit{The Quantum Theory of Fields, Vol.\ I: Foundations},
            Cambridge University Press.
            DOI: 10.1017/CBO9781139644167.
            
            \bibitem{GradshteynRyzhik}
            I.~S.~Gradshteyn and I.~M.~Ryzhik,
            \textit{Table of Integrals, Series, and Products},
            7th ed., Academic Press (2007).
            (Asymptotics and identities for modified Bessel functions.)
            
            \bibitem{PageWootters1983}
            D.~N.~Page and W.~K.~Wootters,
            Evolution without evolution: Dynamics described by stationary observables,
            \textit{Phys.\ Rev.\ D} \textbf{27}, 2885 (1983).
            DOI: 10.1103/PhysRevD.27.2885.
            
            \bibitem{Wootters1984}
            W.~K.~Wootters,
            ``Time'' replaced by quantum correlations,
            \textit{Int.\ J.\ Theor.\ Phys.} \textbf{23}, 701 (1984).
            DOI: 10.1007/BF02214098.
            
            \bibitem{Rovelli2002PartialObservables}
            C.~Rovelli,
            Partial observables,
            \textit{Phys.\ Rev.\ D} \textbf{65}, 124013 (2002).
            DOI: 10.1103/PhysRevD.65.124013.
            
            \bibitem{Dittrich2007PartialComplete}
            B.~Dittrich,
            Partial and complete observables for Hamiltonian constrained systems,
            \textit{Gen.\ Relativ.\ Gravit.} \textbf{39}, 1891 (2007).
            DOI: 10.1007/s10714-007-0495-2.
            
            \bibitem{Tambornino2012RelationalObservables}
            J.~Tambornino,
            Relational observables in gravity: a review,
            \textit{SIGMA} \textbf{8}, 017 (2012).
            arXiv:1109.0740.
            
            \bibitem{Marolf1995AlmostIdealClocks}
            D.~Marolf,
            Almost ideal clocks in quantum cosmology: A brief derivation of time,
            \textit{Class.\ Quantum Grav.} \textbf{12}, 2469 (1995).
            DOI: 10.1088/0264-9381/12/10/007.
            
            \bibitem{GambiniPullin2021}
            R.~Gambini and J.~Pullin,
            Relational time and quantum clocks,
            \textit{Phys.\ Rev.\ D} \textbf{104}, 025003 (2021).
            DOI: 10.1103/PhysRevD.104.025003.
            
            \bibitem{HohnSmith2020}
            P.~A.~H{\"o}hn and A.~R.~H.~Smith,
            Quantum reference frames and time,
            \textit{Phys.\ Rev.\ D} \textbf{101}, 125011 (2020).
            DOI: 10.1103/PhysRevD.101.125011.
            
            \bibitem{Giovannetti2015}
            V.~Giovannetti, S.~Lloyd, and L.~Maccone,
            Quantum time,
            \textit{Nature Physics} \textbf{11}, 1 (2015).
            DOI: 10.1038/nphys3527.
        
        \end{thebibliography}

\end{document}